% TOP_PROBLEM: {"id": 30007041, "title": "A geometric criterion for Kuga–Satake algebraicity", "category_id": 5, "category": {"id": 5, "name": "algebraic_geometry", "display_name": "Algebraic Geometry", "description": "Geometric objects defined by polynomial equations.", "slug": "algebraic-geometry", "order_index": 5, "created_at": "2026-07-31T15:26:25.671Z"}, "status": "open"}
% FIRST_POSTED: 2026-10-06
% ENTRY_KIND: statement_only
% !TeX program = lualatex
\documentclass[11pt]{article}
\usepackage[margin=25mm]{geometry}
\usepackage{fontspec}
\setmainfont{Latin Modern Roman}
\usepackage{amsmath,amssymb,mathtools}
\usepackage{unicode-math}
\setmathfont{Latin Modern Math}
\usepackage{xurl,hyperref}
\hypersetup{colorlinks=true,urlcolor=blue}
\setcounter{secnumdepth}{0}
\setlength{\parindent}{0pt}
\setlength{\parskip}{5pt}
\setlength{\emergencystretch}{3em}
\newcommand{\R}{\mathbb R}
\newcommand{\N}{\mathbb N}
\newcommand{\Z}{\mathbb Z}
\newcommand{\Q}{\mathbb Q}
\newcommand{\C}{\mathbb C}
\newcommand{\D}{\mathbb D}
\newcommand{\F}{\mathbb F}
\newcommand{\sref}[2]{\hyperlink{source:#1:#2}{[S#2]}}
\newcommand{\cataloguescope}[1]{#1}
\begin{document}
\section*{A geometric criterion for Kuga–Satake algebraicity}\label{30007041}
\subsection{Definitions and mathematical statement}
\paragraph{Statement recorded in the supplied source.}
Let (S,h) be a primitively amply polarized projective complex K3 surface. Put $V=h\perp $ in $H^{2}(S,Q)$, with its primitive quadratic Hodge structure, and assume $Hg(V)=SO(V,q)$, so $T(S)=V$. Are the following three conditions equivalent? First, some rational codimension-two correspondence from S to a complex abelian variety acts nontrivially on V. Second, there is a smooth connected projective surface Y and a dominant generically finite morphism $f:Y\to S$ such that tracing the wedge products of holomorphic one-forms on Y to $H^{0}(S,\Omega ^{2}_S)$ gives a nonzero map. Third, for every specified standard full Kuga–Satake Clifford, polarization and realization datum, a rational codimension-two cycle induces the prescribed Kuga–Satake map on $T(S)=V$.
\subsection{Short English statement}
Let (S,h) be a primitively amply polarized projective complex K3 surface. Put $V=h\perp $ in $H^{2}(S,Q)$, with its primitive quadratic Hodge structure, and assume $Hg(V)=SO(V,q)$, so $T(S)=V$. Are the following three conditions equivalent? First, some rational codimension-two correspondence from S to a complex abelian variety acts nontrivially on V. Second, there is a smooth connected projective surface Y and a dominant generically finite morphism $f:Y\to S$ such that tracing the wedge products of holomorphic one-forms on Y to $H^{0}(S,\Omega ^{2}_S)$ gives a nonzero map. Third, for every specified standard full Kuga–Satake Clifford, polarization and realization datum, a rational codimension-two cycle induces the prescribed Kuga–Satake map on $T(S)=V$.
\subsection{Sources}
\begin{itemize}
\item[S1] Original question source or mathematical context.
\url{http://www.openproblemgarden.org/op/the_hodge_conjecture}.
\item[S2] Original question source or mathematical context.
\url{https://arxiv.org/abs/2608.18134}.
\item[S3] Original question source or mathematical context.
\url{https://en.wikipedia.org/wiki/Hodge_conjecture}.
\item[S4] Original question source or mathematical context.
\url{https://www.claymath.org/library/monographs/MPPc.pdf}.
\item[S5] Original question source or mathematical context.
\url{https://www.claymath.org/millennium/hodge-conjecture/}.
\item[S6] Original question source or mathematical context.
\url{https://www.claymath.org/wp-content/uploads/2022/06/hodge.pdf}.
\item[S7] Original question source or mathematical context.
\url{https://zenodo.org/records/22701510}.
\item[S8] OpenAI, manuscript citation page. Source of the formulation motivating this local problem.
\url{https://github.com/openai/math/blob/adc7f1241b42e322a6451854ab7e4b4c146bf78a/preprints/A-Conditional-Reduction-for-Algebraic-Kuga-Satake-Correspondences-September-10-2026/README.md}.
\end{itemize}
\end{document}
